\chapterimage{chapter_pointers.jpg} % Chapter heading image

\chapter{Variáveis, Tipos de Dados e Expressões Aritméticas}

Agora que você já executou seu primeiro programa, é hora de aprender um conceito fundamental: como armazenar e manipular dados. Em todo programa interessante, você precisa guardar informações. Uma idade, um nome, um saldo bancário. C oferece maneiras estruturadas de fazer isso através de variáveis e tipos de dados.

\section{Tipo de dados básicos}

Em C, toda variável deve ter um tipo. O tipo define quanto espaço na memória será reservado e como aquele dado será interpretado.

Os tipos básicos em C são:

\subsection{int (inteiro)}

Representa números inteiros, tanto positivos quanto negativos. Um \textit{int} típico ocupa 4 bytes (32 bits) de memória, permitindo armazenar números de aproximadamente -2 bilhões a +2 bilhões.

\begin{ccode}
  int idade = 25;
  int temperatura = -10;
  int habitantes = 1000000;
\end{ccode}

\subsection{float (ponto flutuante)}

Representa números decimais. Um \textit{float} ocupa 4 bytes e oferece precisão limitada (cerca de 6-7 dígitos significativos).

\begin{ccode}
  float altura = 1.75;
  float pi_aproximado = 3.14;
  float saldo = -50.25;
\end{ccode}

\subsection{double (ponto flutuante duplo)}

Similar ao \textit{float}, mas ocupa 8 bytes e oferece maior precisão (cerca de 15-16 dígitos significativos). Use quando precisar de mais precisão em cálculos matemáticos.

\begin{ccode}
  double pi = 3.141592653589793;
  double resultado_cientifico = 0.000000123456789;
\end{ccode}

\subsection{char (caractere)}

Representa um único caractere. Ocupa 1 byte. Pode ser uma letra, um dígito ou um símbolo especial.

\begin{ccode}
  char letra = 'A';
  char operador = '+';
  char aspas = '"';
\end{ccode}

Lembre-se: caracteres individuais usam aspas simples (\textit{'A'}), enquanto textos completos usam aspas duplas (\textit{"texto"}).

\section{Variáveis}

Uma variável é um local nomeado na memória onde você armazena dados. Antes de usar uma variável, você precisa \textit{declarar} ela.

\begin{quote}
\textit{Declarar? Parece burocrático...}
\end{quote}

É bem simples. Declarar significa dizer ao compilador: \say{Ei, eu preciso de um espaço na memória para guardar um inteiro, e quero chamar esse espaço de 'x'}.

A sintaxe é:

\begin{ccode}
  tipo nome_da_variavel;
\end{ccode}

Exemplos:

\begin{ccode}
  int contador;
  float altura;
  char inicial;
  double resultado;
\end{ccode}

Você também pode \textit{inicializar} uma variável no ato da declaração, ou seja, atribuir um valor inicial:

\begin{ccode}
  int contador = 0;
  float altura = 1.80;
  char inicial = 'J';
  double resultado = 3.14;
\end{ccode}

É uma boa prática sempre inicializar suas variáveis, especialmente em escopos locais. Do contrário, elas conterão lixo (valores aleatórios da memória).

\begin{quote}
\textit{Qual é a diferença entre declaração e inicialização?}
\end{quote}

\textit{Declaração} apenas reserva espaço na memória. \textit{Inicialização} reserva espaço E atribui um valor inicial. Se você declara sem inicializar, o valor será aleatório.

Regras para nomear variáveis:

\begin{itemize}
  \item Começam com uma letra (a-z, A-Z) ou sublinhado (\_)
  \item Podem conter letras, dígitos (0-9) e sublinhadosDepois de caracteres iniciais
  \item São \textit{case-sensitive}, ou seja, \textbf{contador} é diferente de \textbf{Contador}
  \item Não podem ter espaços
  \item Não podem ser palavras reservadas (like \textit{int}, \textit{while}, \textit{if}, etc)
\end{itemize}

Bons nomes de variáveis:

\begin{ccode}
  int idade_usuario;
  float saldo_conta;
  char inicial_nome;
  int contador_loops;
\end{ccode}

Nomes ruins (evite):

\begin{ccode}
  int x;           // Muito vago
  int DadosDoSistema;  // Muito genérico
  int a1b2c3;      // Ilegível
\end{ccode}

\section{Modificadores}

Modificadores alteram o comportamento dos tipos básicos. Os principais são:

\subsection{unsigned (sem sinal)}

Por padrão, tipos inteiros podem representar números negativos. Se você sabe que seus dados são sempre positivos ou zero, use \textit{unsigned} para dobrar o intervalo positivo.

\begin{ccode}
  int idade = -5;              // Permitido
  unsigned int quantidade = 1000;  // Sempre positivo
  unsigned int quantidade_negativa = -5;  // ERRO! Undefined behavior
\end{ccode}

\subsection{short (curto) e long (longo)}

Modificam o tamanho do tipo:

\begin{ccode}
  short int valor_pequeno;     // Geralmente 2 bytes
  long int valor_grande;       // Geralmente 8 bytes
  long double precisao_dupla;  // Maior precisão que double
\end{ccode}

\subsection{signed (com sinal)}

Explicitamente permite números negativos (é o padrão para \textit{int} e \textit{char}).

\begin{ccode}
  signed int temperatura;      // Pode ser negativo
\end{ccode}

\section{Constantes}

Uma constante é um valor que não pode ser alterado após sua inicialização. Use a palavra-chave \textit{const}:

\begin{ccode}
  const float pi = 3.14159;
  const int dias_semana = 7;
  const char separador = '/';
\end{ccode}

Se você tentar modificar uma constante, o compilador dará um erro:

\begin{ccode}
  const int velocidade_luz = 299792458;
  velocidade_luz = 300000;  // ERRO!
\end{ccode}

Constantes melhoram a legibilidade e evitam erros acidentais. Se um valor não deve ser alterado, declareNELE como constante.

\section{Expressões Aritméticas}

Você pode realizar operações matemáticas entre variáveis. Os operadores básicos são:

\subsection{Adição (+)}

\begin{ccode}
  int a = 5;
  int b = 3;
  int soma = a + b;  // soma = 8
\end{ccode}

\subsection{Subtração (-)}

\begin{ccode}
  int x = 10;
  int y = 3;
  int diferenca = x - y;  // diferenca = 7
\end{ccode}

\subsection{Multiplicação (*)}

\begin{ccode}
  int multiplicador = 4;
  int multiplicando = 5;
  int produto = multiplicador * multiplicando;  // produto = 20
\end{ccode}

\subsection{Divisão (/)}

\begin{quote}
\textit{Cuidado!}
\end{quote}

A divisão entre dois inteiros resulta em um inteiro. Qualquer parte decimal é perdida:

\begin{ccode}
  int dividendo = 7;
  int divisor = 2;
  int quociente = dividendo / divisor;  // quociente = 3 (não 3.5!)
\end{ccode}

Para obter resultado decimal, use \textit{float} ou \textit{double}:

\begin{ccode}
  float resultado = 7.0 / 2.0;  // resultado = 3.5
\end{ccode}

\subsection{Resto (\%)}

O operador de módulo retorna o resto de uma divisão:

\begin{ccode}
  int dividendo = 17;
  int divisor = 5;
  int resto = dividendo \% divisor;  // resto = 2 (pois 17 = 5*3 + 2)
\end{ccode}

Este operador é frequentemente usado para verificar se um número é par ou ímpar:

\begin{ccode}
  int numero = 10;
  if (numero \% 2 == 0) {
    // numero é par
  } else {
    // numero é ímpar
  }
\end{ccode}

\subsection{Precedência de Operadores}

Assim como na matemática, existe uma ordem de precedência. Multiplicação e divisão ocorrem antes de adição e subtração:

\begin{ccode}
  int resultado = 2 + 3 * 4;  // resultado = 14, não 20
  // Primeiro: 3 * 4 = 12
  // Depois: 2 + 12 = 14
\end{ccode}

Use parênteses para deixar claro a ordem:

\begin{ccode}
  int resultado1 = (2 + 3) * 4;  // resultado1 = 20
  int resultado2 = 2 + (3 * 4);  // resultado2 = 14
\end{ccode}

\section{Conversão de tipo de dados}

Às vezes você precisa converter um tipo de dado para outro. C oferece duas formas: conversão implícita e conversão explícita.

\subsection{Conversão Implícita}

O compilador automaticamente converte tipos quando faz sentido:

\begin{ccode}
  int x = 5;
  float y = x;  // int é automaticamente convertido para float
  // y agora é 5.0
\end{ccode}

\subsection{Conversão Explícita (Cast)}

Você pode forçar uma conversão usando a sintaxe \textit{(tipo)valor}:

\begin{ccode}
  float pi = 3.14159;
  int inteiro = (int)pi;  // inteiro = 3 (parte decimal descartada)

  int velocidade = 100;
  float velocidade_decimal = (float)velocidade;  // 100.0
\end{ccode}

\begin{quote}
\textit{Cuidado ao converter!}
\end{quote}

Conversões podem perder dados. Converter de \textit{float} para \textit{int} descarta a parte decimal. Converter um número grande para um tipo menor pode causar estouro.

Exemplo prático:

\begin{ccode}
  #include <stdio.h>

  int main(void)
  {
    int nota1 = 85;
    int nota2 = 90;
    int nota3 = 78;
    float media = (nota1 + nota2 + nota3) / 3.0;

    printf("Média: %.2f\n", media);
    return 0;
  }
\end{ccode}

Este programa calcula a média de três notas. Note o \textit{3.0} (em vez de \textit{3}) para garantir divisão em ponto flutuante.

Agora você entende os fundamentos de variáveis e tipos de dados em C. No próximo capítulo, aprenderemos como controlar o fluxo do programa usando decisões (if, else, switch).

%------------------------------------------------



